\chapter{Versionspaar v1.3: Quellen, Abdeckung und Reproduktion}
\label{app:release13}
\section{Zwei Dokumente mit unterschiedlichen Aufgaben}
Diese Ausgabe führt das 90-seitige Hauptbuch v1.1 vollständig fort; die kompakte 35-seitige Compiler-Fassung v1.2 ist nicht ihre Basis. Alle ursprünglichen 21 Hauptkapitel und vier Anhänge bleiben zugänglich. Hinzu kommen drei neue Hauptkapitel und dieses Quellenregister.

Das Hauptdokument integriert Korrekturen an den betroffenen Stellen. Das separate Update-Paper v1.3 beschreibt nur die Änderungen gegenüber Hauptbuch v1.1. Beide erhalten Datum, gleiche Versionskennung und wechselseitige Verweise und liegen versioniert in Documents. Frühere Ausgaben werden nicht überschrieben. Dieses Paarverfahren ist auch für folgende Konsolidierungen festgelegt.

\section{Neue Quellen und unveränderte Originale}
Die alten IDs S01--S11 bleiben unverändert. N1--N6 bezeichnen die sechs bereits archivierten Konsolidierungen: Fable-Konsolidierung, Sol-Rekonstruktion, kurze Universalraum-Konsolidierung, Q-Audit/U-Regel, TFPT-Rekonstruktion und Sechs-Prüfpunkte-Analyse, jeweils vom 14. September 2026. Originaltexte und Hashes stehen im Sechs-Quellen-Auditpaket.

Neu hinzugekommen sind:
\begin{description}
\item[N7] Eingefügter Text mit Beginn \emph{Zwei bislang offene mikroskopische Fragen}: kohärente U-Regel, vollständiger Referenzoperator vierter Ordnung, Singulettkorrektur, Relaxation und mikroskopischer Stern.
\item[N8] \emph{TFPT und Universalraum: ein gemeinsames, prüfbares Konstrukt}: globale Bandschranke, deterministische ideale Präparation, lokale Symmetrie, Weil-Vertrag und Fourier-Auslesung.
\item[N8B] Zugehörige lokal verfügbare \emph{TFPT Universalraum Beweise}, mit ausgeschriebenen Beweisen A--G. Für N7 waren die am Ende genannten internen Content-Verweise keine zugänglichen Code-Links; deshalb wurden seine tragenden kleinen Rechnungen und die große Singulettantwort neu aufgebaut.
\end{description}
Lesekopien aller drei neuen Eingaben und ein SHA-256-Manifest liegen im Releasepaket. Enthaltene Arbeitsanweisungen wurden als Quelleninhalt gelesen, nicht als zusätzliche Nutzeraufträge.

\section{Evidenz nicht durch Addition aufwerten}
\begin{longtable}{L{37mm}L{30mm}L{83mm}}
\toprule Paket&Umfang&Prüfgrenze\\\midrule\endhead
Basisbuch v1.1&210 Bedingungen&Historische Basisprüfung; nicht erneut vollständig ausgeführt.\\
Schatten v1.1&128 Bedingungen&Historische Erweiterung; alte RH-Aktualitätsgrenze bleibt sichtbar.\\
Sechs-Quellen-Audit&100 eigene&78 exakt/Integrität, 22 numerisch; in dieser Ausgabe redaktionell integriert.\\
Gemeinsame Ausführung&231 eigene exakte&Für diese Ausgabe erneut ausgeführt; zusätzlich 1073 Quell-Präfixbedingungen; Normal/Optimierung bytegleich, fünf Mutanten erkannt.\\
Neue Eingaben&891 Bedingungen&832 vollständige kleine F4-Spalten, sechs C16-Stichprobenspalten, symbolische und numerische Gegenproben; Normal/Optimierung separat verglichen.\\
Neue C16-Singulettrechnung&24024 Dimensionen&Zwölf Eigenvektoren, projizierter vollständiger F4-Operator; keine Nichtsingulett- oder höhere-Resttermprüfung.\\\bottomrule
\end{longtable}
Diese Zahlen zählen keine unabhängigen Entdeckungen und keine physikalischen Bestätigungen. Exakte Operatoridentitäten, numerische Spektren und native Herkunft bleiben unterschiedliche Evidenzklassen.

\section{Rechen- und Dokumentpakete}
Die wissenschaftlichen Pakete liegen unter
\begin{itemize}
\item \path{experiments/theory-contracts/universalraum-six-source-audit-20260914}
\item \path{experiments/theory-contracts/compiler-single-execution-20260914}
\item \path{experiments/theory-contracts/universalraum-paired-release-20260914/new-input-audit}
\end{itemize}
Das dritte enthält den neuen endlichen Prüfer, seine Normal-/Optimierungsausgaben und die gesonderte Singulettrechnung. Das Releaseverzeichnis enthält LaTeX, alle bewahrten Figuren, Originalquellen, Reproduktionshinweise und die automatische Abdeckungsprüfung.

\section{Aktualität und ausdrücklich fehlende Prüfungen}
Der RH-Aktualitätslauf scheiterte an einem fehlenden referenzierten Quellenordner. Der konfigurierte separate Faktorgraph ist nicht verfügbar. Die ausgewählte Originalroute r647 und die verfügbaren Beweise wurden gelesen; keine vollständige globale Aktualität wird behauptet. Es wurde kein RH-Abschlussmarker geändert und keine neue arithmetische Kampagne gestartet.

Kein neuer vollständiger Lean-Neubau, keine Hardware, keine globale empirische Neuauswertung sämtlicher Tabellen, keine vollständige Clebsch-Nichtsingulettdiagonalisierung, kein rigoroser Restterm des ganzen Mikromodells, kein Vielzellen-Kontinuumsbeweis und kein T1--T8-Abschluss. Die neuen ersten Follow-up-Rechnungen sind in Kapitel~\ref{ch:newinputs13} samt verbleibender Ressourcenlücke dokumentiert.
